\documentclass[11pt]{article}

\usepackage[letterpaper,margin=1in]{geometry}
\usepackage{amsmath,amssymb,amsthm}
\usepackage{booktabs,array}
\usepackage{xcolor}
\usepackage{microtype}
\usepackage[colorlinks=true,allcolors=blue!55!black]{hyperref}

\newtheorem{theorem}{Theorem}[section]
\newtheorem{lemma}[theorem]{Lemma}
\newtheorem{proposition}[theorem]{Proposition}
\newtheorem{remark}[theorem]{Remark}

\newcommand{\E}{\mathbb{E}}
\newcommand{\V}{\operatorname{Var}}
\newcommand{\N}{\mathcal{N}}
\newcommand{\Bin}{\operatorname{Binomial}}
\newcommand{\Pois}{\operatorname{Poisson}}
\newcommand{\Ber}{\operatorname{Bernoulli}}
\newcommand{\logit}{\operatorname{logit}}
\newcommand{\sigmoid}{\operatorname{logit}^{-1}}
\newcommand{\softplus}{\operatorname{softplus}}
\newcommand{\KL}{\operatorname{KL}}
\newcommand{\ELBO}{\operatorname{ELBO}}
\newcommand{\R}{\mathbb{R}}

\title{A splitting approach for variational inference in\\
non-Gaussian latent Gaussian models}
\author{Dongyue Xie \and Peter Carbonetto \and Matthew Stephens}
\date{Revised working draft: \today}

\begin{document}
\maketitle

\begin{abstract}
Many analyses of high-dimensional biological data begin by transforming
non-Gaussian observations and then applying a Gaussian method. Examples
include log-transforming RNA sequencing counts before factor analysis and
transforming accessibility proportions before functional regression. These
transformations are convenient, but they discard information about the
observation process and can be sensitive to arbitrary choices such as a
pseudocount. We develop a modular variational framework that couples an
exponential-family observation model to an existing empirical Bayes Gaussian
method through a latent Gaussian splitting variable. The resulting algorithm
alternates between independent low-dimensional variational updates for the
non-Gaussian observations and a Gaussian update that can be delegated to
well-developed software for normal means, regression, smoothing, matrix
factorization, multivariate analysis, or structured variable selection.

We give explicit updates for Poisson and binomial observations. The binomial
case is especially useful when a response is an aggregate of binary
measurements: the number of successes and the number of trials are retained,
so observations with different sampling depths contribute different amounts
of information. We also clarify an important limitation. With one Bernoulli
observation per latent mean, the splitting variance can be weakly identified
or non-identifiable; replication within a binomial observation supplies the
extra-binomial information needed to estimate this variance. For pure
Bernoulli data we describe fixed-variance and data-augmentation alternatives.
The framework therefore provides both an algorithm and a diagnostic guide for
deciding when split variational inference is appropriate. We conclude with a
binomial functional regression model for single-molecule chromatin
accessibility data, in which wavelet-based functional fine-mapping is used as
the Gaussian module on the latent log-odds scale.
\end{abstract}

\section{Introduction}

Consider observations $y_i$ whose distribution belongs to an exponential
family,
\begin{equation}
  y_i\mid \mu_i \sim \mathcal{D}\{h(\mu_i)\},
  \qquad
  \mu=(\mu_1,\ldots,\mu_m)^\mathsf{T}\sim f,
  \label{eq:general-model}
\end{equation}
where $\mu_i$ is a canonical or transformed mean parameter and $h$ is the
inverse link. Depending on the structure assigned to $f$, this formulation
includes empirical Bayes normal means, sparse regression, smoothing,
multivariate shrinkage, and matrix factorization. The corresponding Gaussian
problems are often computationally mature, whereas replacing Gaussian
observations by counts or proportions may require a new inference algorithm
for every choice of $f$.

Our goal is to reuse these Gaussian methods without first transforming the
data. We introduce a latent Gaussian layer that separates the likelihood from
the structured prior. This device is analogous in spirit to splitting methods
in optimization: a difficult coupled problem is replaced by two simpler
problems connected through low-dimensional posterior moments. One module
handles the non-Gaussian likelihood; the other calls an empirical Bayes
Gaussian method (EBG).

The original motivation was dimension reduction for single-cell RNA
sequencing counts. A second important use arises when binary measurements are
aggregated within a biological sample. In Fiber-seq, for example, accessibility
is observed on individual chromatin fibers but genetic covariates are measured
at the donor level. The statistically relevant observation is then a number of
accessible fibers out of a known number assayed. An analysis of the proportion
alone cannot distinguish, say, one accessible fiber out of two from one
hundred accessible fibers out of two hundred. A binomial likelihood preserves
this distinction and avoids treating fibers from the same donor as independent
genetic samples.

Table~\ref{tab:gaussian-modules} lists examples of structured priors and the
Gaussian module they induce.

\begin{table}[ht]
\centering
\caption{Examples of structured priors and reusable Gaussian modules.}
\label{tab:gaussian-modules}
\begin{tabular}{>{\raggedright\arraybackslash}p{0.33\linewidth}
                >{\raggedright\arraybackslash}p{0.28\linewidth}
                >{\raggedright\arraybackslash}p{0.27\linewidth}}
\toprule
Structure for $b$ & Model & Example Gaussian engine \\
\midrule
$b_i\overset{\mathrm{iid}}\sim g$ & univariate mean inference & \texttt{ash}, \texttt{ebnm} \\
$b=X\beta$, $\beta\sim g$ & regression & sparse EB regression \\
$b$ spatially smooth & smoothing & wavelet or Gaussian-process EB \\
$b=LF^\mathsf{T}$ & matrix factorization & \texttt{flash} \\
$b_i\sim\sum_k\pi_k\N(0,U_k)$ & multivariate mean inference & \texttt{mash} \\
$b=XB$, structured rows of $B$ & functional variable selection & fSuSiE \\
\bottomrule
\end{tabular}
\end{table}

\section{Split variational inference}

\subsection{The augmented model}

Introduce a splitting variable $b=(b_1,\ldots,b_m)^\mathsf{T}$ and write
\begin{align}
 y_i\mid\mu_i &\sim \mathcal{D}\{h(\mu_i)\}, \nonumber\\
 \mu_i\mid b_i,\tau_i^2 &\sim \N(b_i,\tau_i^2),
 \label{eq:split-model}\\
 b &\sim g. \nonumber
\end{align}
The variance $\tau_i^2$ can be interpreted in two ways. It may be an
algorithmic device that connects the likelihood to a Gaussian engine, or it
may be part of the scientific model, representing overdispersion or residual
variation on the latent scale. This distinction matters for identifiability
and interpretation.

We use a mean-field approximation
\[
 q(\mu,b)=q_\mu(\mu)q_b(b)
 =\left\{\prod_i q_{\mu_i}(\mu_i)\right\}q_b(b).
\]
Up to constants that do not depend on $q$, the evidence lower bound is
\begin{align}
 \mathcal{F}(q,g;\tau^2)
  ={}&\sum_i\E_{q_{\mu_i}}\log p(y_i\mid\mu_i)
   +\sum_i\E_q\log p(\mu_i\mid b_i,\tau_i^2) \nonumber\\
  &+\E_{q_b}\log\frac{g(b)}{q_b(b)}
   -\sum_i\E_{q_{\mu_i}}\log q_{\mu_i}(\mu_i).
 \label{eq:split-elbo}
\end{align}
We take $q_{\mu_i}=\N(m_i,v_i)$. Let
$\bar b_i=\E_{q_b}b_i$ and $V_{b_i}=\V_{q_b}(b_i)$. Conditional on the
current Gaussian module, the update of $(m_i,v_i)$ is independent across
$i$ and depends on $q_b$ only through $\bar b_i$; $V_{b_i}$ contributes a
constant to this coordinate update.

\subsection{The modular algorithm}

For a common splitting variance $\tau^2$, coordinate ascent alternates among
three steps:

\begin{enumerate}
\item \textbf{Non-Gaussian step.} For each $i$, update
  $q_{\mu_i}=\N(m_i,v_i)$ using the likelihood for $y_i$, prior mean
  $\bar b_i$, and prior variance $\tau^2$.
\item \textbf{Gaussian step.} Treat $m=(m_1,\ldots,m_m)^\mathsf{T}$ as a
  Gaussian response with residual variance $\tau^2$, and update $q_b$ and
  any hyperparameters in $g$ with the selected EBG method.
\item \textbf{Variance step.} If $\tau^2$ is estimated, update
\begin{equation}
 \tau^2\leftarrow \frac{1}{m}\sum_i
 \left\{(m_i-\bar b_i)^2+v_i+V_{b_i}\right\}.
 \label{eq:tau-update}
\end{equation}
\end{enumerate}

The posterior variance $V_{b_i}$ in~\eqref{eq:tau-update} is essential. Omitting
it generally underestimates the latent residual variance. The same pattern
extends to matrices: average the corresponding quantity over all entries.
When $\tau_i^2$ is known and heterogeneous, the Gaussian module must support
heterogeneous observation variances.

\subsection{Relation to the collapsed objective}

Let
\[
 f_\tau(\mu)=\int \phi_\tau(\mu-b)\,dG(b)
\]
be the Gaussian convolution of $G$, where $\phi_\tau$ is a normal density with
variance $\tau^2$.

\begin{lemma}[Curvature of a Gaussian convolution]
For every distribution $G$ for which the convolution is finite,
\begin{equation}
 \frac{d^2}{d\mu^2}\log f_\tau(\mu)
 =-\frac{1}{\tau^2}
  +\frac{\V(b\mid\mu)}{\tau^4}
 \geq -\frac{1}{\tau^2}.
 \label{eq:curvature}
\end{equation}
\end{lemma}

\begin{proof}
Differentiate under the integral. The first derivative is
$(\E[b\mid\mu]-\mu)/\tau^2$; differentiating once more yields
\eqref{eq:curvature}.
\end{proof}

Suppose $q_\mu=\N(m,v)$. The curvature bound implies
\begin{equation}
 \E_{q_\mu}\log f_\tau(\mu)
 \geq \log f_\tau(m)-\frac{v}{2\tau^2}.
 \label{eq:collapsed-bound}
\end{equation}
The right-hand side equals the result of profiling the splitting ELBO over
$q_b$. Thus the split objective is a lower bound on the variational objective
obtained after integrating out $b$. Equality holds when $G$ is a point mass;
with a flexible $G$, the gap is the price paid for modularity.

\section{Variational Gaussian updates}

\subsection{Poisson observations}

Let
\[
 y_i\mid\mu_i\sim\Pois\{s_i\exp(\mu_i)\},
 \qquad s_i>0,
\]
and let the current prior mean be $\bar b_i$. Up to constants, the local
objective for $q_{\mu_i}=\N(m_i,v_i)$ is
\begin{align}
 \mathcal{F}_i(m_i,v_i)
  ={}&y_i m_i-s_i\exp(m_i+v_i/2)
  -\frac{(m_i-\bar b_i)^2+v_i}{2\tau^2}
  +\frac{1}{2}\log v_i.
 \label{eq:poisson-vga}
\end{align}
The negative of~\eqref{eq:poisson-vga} is strictly convex. The stationary
conditions are
\begin{align}
 0&=y_i-s_i\exp(m_i+v_i/2)-\frac{m_i-\bar b_i}{\tau^2},
 \label{eq:pois-score-m}\\
 \frac{1}{v_i}&=s_i\exp(m_i+v_i/2)+\frac{1}{\tau^2}.
 \label{eq:pois-score-v}
\end{align}
They can be solved by damped Newton iterations. Alternatively, substituting
one equation into the other reduces the update to a stable one-dimensional
root-finding problem. Equation~\eqref{eq:pois-score-v} also shows that
$0<v_i<\tau^2$.

\subsection{Binomial observations}

Let
\begin{equation}
 k_i\mid\mu_i\sim\Bin\{n_i,\sigmoid(\mu_i)\},
 \qquad 0\leq k_i\leq n_i,
 \label{eq:binomial-model}
\end{equation}
where $k_i$ is the number of successes and $n_i$ is the known number of
trials. The local objective is
\begin{align}
 \mathcal{F}_i(m_i,v_i)
 ={}& k_i m_i-n_i A(m_i,v_i)
 -\frac{(m_i-\bar b_i)^2+v_i}{2\tau^2}
 +\frac{1}{2}\log v_i,
 \label{eq:binomial-vga}\\
 A(m,v)={}&\E_{Z\sim\N(m,v)}\log(1+e^Z).
 \nonumber
\end{align}
There is no elementary expression for $A(m,v)$, but Gauss-Hermite quadrature
is accurate and inexpensive:
\begin{equation}
 A(m,v)\approx \frac{1}{\sqrt{\pi}}
 \sum_{r=1}^R w_r\log\left\{1+
 \exp\left(m+\sqrt{2v}\,z_r\right)\right\},
 \label{eq:gh}
\end{equation}
where $(z_r,w_r)$ are fixed Hermite nodes and weights. Stable implementations
evaluate the logarithm with a softplus function.

Useful exact derivative identities are
\begin{align}
 \frac{\partial A}{\partial m}
 &=\E\{p(Z)\}, \\
 \frac{\partial A}{\partial v}
 &=\frac{1}{2}\E\{p(Z)[1-p(Z)]\},
 \qquad p(z)=\sigmoid(z).
\end{align}
The local first-order conditions therefore become
\begin{align}
 0={}&k_i-n_i\E\{p(Z_i)\}
      -\frac{m_i-\bar b_i}{\tau^2},
 \label{eq:bin-score-m}\\
 \frac{1}{v_i}={}&\frac{1}{\tau^2}
 +n_i\E\{p(Z_i)[1-p(Z_i)]\}.
 \label{eq:bin-score-v}
\end{align}
As in the Poisson case, $0<v_i<\tau^2$. We recommend optimizing in
$(m_i,\log v_i)$ with a trust-region or damped Newton method, using the same
quadrature rule for the objective and derivatives. Warm starts from the
previous outer iteration make these scalar updates inexpensive.

\subsection{Bernoulli identifiability and the role of replication}

The binomial extension requires more care than the Poisson extension because
the splitting variance need not be identified by Bernoulli data.

\begin{proposition}[A Bernoulli non-identifiability case]
Suppose $x_i\mid\mu_i\sim\Ber\{\sigmoid(\mu_i)\}$ and
$\mu_i\sim\N(0,\tau^2)$. Then
$\Pr(x_i=1)=1/2$ for every $\tau^2>0$; hence the Bernoulli likelihood contains
no information about $\tau^2$.
\end{proposition}

\begin{proof}
The normal distribution is symmetric around zero and
$\sigmoid(z)+\sigmoid(-z)=1$. Therefore
$\E\{\sigmoid(\mu_i)\}=1/2$ for all $\tau^2$.
\end{proof}

With nonzero or varying prior means, a Bernoulli likelihood may contain some
information about $\tau^2$, but it is often weak. Moreover, if the structured
prior $g$ has an estimable scale, variation can be shifted between $g$ and
$\tau^2$, creating a ridge or boundary solution.

Replication within a binomial observation changes this. If
$K\mid P\sim\Bin(n,P)$, then
\begin{equation}
 \V(K)=n\bar p(1-\bar p)+n(n-1)\V(P),
 \qquad \bar p=\E(P).
 \label{eq:binomial-overdispersion}
\end{equation}
The second term vanishes when $n=1$ and grows quadratically with $n$.
Consequently, aggregated binomial data can identify latent heterogeneity from
extra-binomial variation, whereas a single Bernoulli trial cannot. Practical
identifiability still depends on the range of $n_i$, the number of independent
biological units, the range of fitted probabilities, and the flexibility of
$g$.

We recommend three safeguards:
\begin{enumerate}
\item retain the trial totals $n_i$ and diagnose their distribution rather
than analyzing proportions alone;
\item compare estimated and fixed values of $\tau^2$, and monitor the ELBO for
flat profiles or boundary optima;
\item use null simulations or permutations that preserve the observed trial
totals and dependence structure.
\end{enumerate}

\subsection{Alternative augmentations for binary data}

When all $n_i=1$, a fixed splitting variance may still provide a useful
computational approximation, but it should not be interpreted as an estimated
overdispersion parameter. Two classic alternatives are worth considering.
Albert--Chib augmentation introduces a truncated latent normal and is exact
for a \emph{probit} likelihood, $\Pr(x_i=1)=\Phi(b_i)$; it is not an exact
logistic representation. Polya--Gamma augmentation is conditionally Gaussian
for logistic binomial models and provides exact augmentation before the
mean-field approximation. Its Gaussian pseudo-observations have
observation-specific precisions, so the downstream Gaussian engine must accept
heterogeneous weights. Split VA trades this extra Gaussian-engine complexity
for independent two-parameter local updates and a common latent residual
variance.

\section{Reusable Gaussian modules}

\subsection{Empirical Bayes normal means}

If $b_i\overset{\mathrm{iid}}\sim g$, the Gaussian step solves
\[
 m_i\mid b_i\sim\N(b_i,\tau^2),\qquad b_i\sim g.
\]
This is the empirical Bayes normal means problem. Scale mixtures of normals,
unimodal families, or nonparametric priors can be fitted with existing
software. The returned posterior means and variances feed directly into the
next non-Gaussian and variance updates.

\subsection{Regression and smoothing}

Let $b=X\beta$ or let $b$ be a smooth function evaluated on a grid. The
Gaussian step is ordinary empirical Bayes regression or smoothing with
response $m$ and residual variance $\tau^2$. The prior on $\beta$ can induce
sparsity, and the prior on a smooth function can be a Gaussian process or a
sparse prior on wavelet coefficients. This modularity allows the same Poisson
or binomial likelihood code to be paired with several notions of structure.

\subsection{Matrix factorization}

For observations $y_{ij}$, let
\begin{align}
 y_{ij}\mid\mu_{ij}&\sim\mathcal{D}\{h(\mu_{ij})\},\\
 \mu_{ij}\mid L,F&\sim
 \N\left(\sum_{r=1}^K l_{ir}f_{jr},\tau_{ij}^2\right).
\end{align}
A mean-field approximation
$q_\mu\prod_r q(l_r)q(f_r)$ yields a non-Gaussian update for every matrix
entry and a Gaussian factor-analysis update using the matrix of latent means
$M=(m_{ij})$. Sparsity, non-negativity, or smoothness can be imposed through
the factor-analysis priors without changing the observation module.

\subsection{Multivariate mean inference}

For multivariate effects, the Gaussian step can use a mixture
\[
 b_i\sim\sum_k\pi_k\N(0,U_k).
\]
The non-Gaussian module returns a vector of latent posterior means for every
unit, and a multivariate adaptive-shrinkage method updates the mixture weights
and posterior moments. Heterogeneous splitting variances require the
multivariate engine to accept observation-specific standard errors.

\subsection{Binomial functional regression and Fiber-seq}
\label{sec:fiberseq}

Let $A_{irt}\in\{0,1\}$ record whether fiber $r$ from donor $i$ is accessible
at genomic position $t$. Define
\[
 K_{it}=\sum_{r=1}^{N_{it}}A_{irt},
 \]
where $N_{it}$ is the number of fibers observed for donor $i$ at position $t$.
The proposed model is
\begin{align}
 K_{it}\mid\mu_{it}
 &\sim\Bin\{N_{it},\sigmoid(\mu_{it})\},
 \label{eq:fiber-binomial}\\
 \mu_{it}\mid a_t,B,\tau^2
 &\sim\N\{a_t+(XB)_{it},\tau^2\}.
 \label{eq:fiber-latent}
\end{align}
Here $X$ contains donor-level genotypes and covariates, $a_t$ is a baseline
accessibility curve, and the rows of $B$ are SNP effect curves on the log-odds
scale. A functional Sum of Single Effects prior can be placed on $B$: each
single effect selects one SNP, while a sparse wavelet prior regularizes its
effect curve across positions.

This model preserves three important features of the experiment. First, the
binomial likelihood distinguishes different coverage totals. Second, donors,
not fibers, remain the independent units for genetic association. Third, the
wavelet transform is applied to latent effect functions rather than directly
to binary observations, so the mean--variance relationship is handled by the
likelihood.

One iteration of binomial fSuSiE is:
\begin{enumerate}
\item update every $q(\mu_{it})=\N(m_{it},v_{it})$ using
\eqref{eq:binomial-vga}, with prior mean
$\bar a_t+\E_q(XB)_{it}$ and variance $\tau^2$;
\item fit Gaussian fSuSiE to $M=(m_{it})$, holding its residual variance fixed
at the current $\tau^2$;
\item update the baseline curve and compute the first two posterior moments of
$a_t+(XB)_{it}$;
\item update $\tau^2$ by the matrix version of~\eqref{eq:tau-update}.
\end{enumerate}

The raw inputs should therefore contain both $K_{it}$ and $N_{it}$. Supplying
only $K_{it}/N_{it}$ loses sampling precision. If fibers are more correlated
than model~\eqref{eq:fiber-latent} can accommodate, one can enrich the latent
model with donor- or region-specific random effects. Inference should be
validated with simulations that reproduce the observed coverage distribution,
baseline accessibility, within-fiber spatial dependence, and linkage
disequilibrium among SNPs.

\section{Computation and diagnostics}

The local Poisson and binomial updates are embarrassingly parallel. Their cost
is linear in the number of observations times the number of quadrature nodes;
the Gaussian module usually dominates total running time. Useful numerical and
statistical diagnostics include:
\begin{itemize}
\item monotonicity of the full split ELBO, evaluated with the same quadrature
rule used by the updates;
\item relative changes in latent means, fitted structured means, and
$\tau^2$;
\item the frequency of local solver failures and boundary values of $v_i$;
\item profile plots of the objective over $\tau^2$;
\item calibration under null simulations or covariate permutations;
\item sensitivity to quadrature order, initialization, and fixed versus
estimated splitting variance.
\end{itemize}

For binomial functional regression, additional diagnostics should compare
observed and fitted success counts across coverage strata and positions. A
method that fits mean accessibility but fails to reproduce the variance across
donors is not adequately modeling the data, even if its fitted mean curves
look plausible.

\section{Discussion}

Split variational inference makes non-Gaussian modeling modular: the
likelihood is handled by simple local variational problems, while structural
learning is delegated to a Gaussian method. This division is particularly
attractive when the Gaussian method is sophisticated, such as adaptive
shrinkage, multivariate empirical Bayes, matrix factorization, or functional
fine-mapping.

The binomial case adds both an important application and an important warning.
Aggregated binary data with known trial totals are not merely ``continuous
values between zero and one.'' They are binomial counts whose precision
depends on the number of trials and the latent success probability. Split VA
naturally uses this information. By contrast, pure Bernoulli data can provide
little or no information about an extra latent variance, particularly near a
zero linear predictor. In that setting the variance should be fixed or
regularized, or a direct probit or Polya--Gamma augmentation should be used.

For Fiber-seq fine-mapping, the binomial formulation is scientifically more
defensible than applying a homoskedastic Gaussian model to mean accessibility.
It is also computationally realistic because the current Gaussian fSuSiE
engine can be reused once it accepts a latent response and a fixed residual
variance. The main research tasks are therefore the binomial local solver,
careful propagation of posterior uncertainty, identifiability diagnostics,
and realistic calibration experiments rather than a redesign of functional
fine-mapping itself.

\begin{thebibliography}{99}

\bibitem{albertchib1993}
Albert, J. H. and Chib, S. (1993).
Bayesian analysis of binary and polychotomous response data.
\emph{Journal of the American Statistical Association}, 88:669--679.

\bibitem{blei2017}
Blei, D. M., Kucukelbir, A., and McAuliffe, J. D. (2017).
Variational inference: A review for statisticians.
\emph{Journal of the American Statistical Association}, 112:859--877.

\bibitem{boyd2011}
Boyd, S., Parikh, N., Chu, E., Peleato, B., and Eckstein, J. (2011).
Distributed optimization and statistical learning via the alternating
direction method of multipliers.
\emph{Foundations and Trends in Machine Learning}, 3:1--122.

\bibitem{denault2025}
Denault, W. R. P. et al. (2025).
fSuSiE enables fine-mapping of QTLs from genome-scale molecular profiles.
\emph{bioRxiv}, doi:10.1101/2025.08.17.670732.

\bibitem{jaakkola1997}
Jaakkola, T. S. and Jordan, M. I. (1997).
A variational approach to Bayesian logistic regression models and their
extensions.
In \emph{Proceedings of the Sixth International Workshop on Artificial
Intelligence and Statistics}, 283--294.

\bibitem{polson2013}
Polson, N. G., Scott, J. G., and Windle, J. (2013).
Bayesian inference for logistic models using Polya--Gamma latent variables.
\emph{Journal of the American Statistical Association}, 108:1339--1349.

\bibitem{stephens2017}
Stephens, M. (2017).
False discovery rates: a new deal.
\emph{Biostatistics}, 18:275--294.

\bibitem{wang2020}
Wang, G., Sarkar, A., Carbonetto, P., and Stephens, M. (2020).
A simple new approach to variable selection in regression, with application
to genetic fine mapping.
\emph{Journal of the Royal Statistical Society: Series B}, 82:1273--1300.

\end{thebibliography}

\end{document}
